\ficha{Flandreau e Zumer (2004), The Making of Global Finance}{flandreauzumer2004}

\begin{identificacao}
FLANDREAU, Marc; ZUMER, Frédéric. \emph{The Making of Global Finance
1880-1913}. Paris: OECD Development Centre, 2004. (Development Centre
Studies). DOI: 10.1787/9789264015364-en.

Monografia de história econômica quantitativa, em inglês, 148 páginas no PDF.
Leitura seletiva: Introdução (p. 11-16), seções I a III (p. 17-26), VI a VIII
(p. 36-47) e X (p. 55-58), com as notas correspondentes (p. 59-70). Consulta
pontual ao fim da seção V (p. 33-35), à seção IX (p. 48-53) para os testes
``trade'' e ``tax'', ao Apêndice Estatístico (p. 96-104) e às tabelas da base
(p. 125-131). Não lidos: seções IV, V e IX na íntegra, e o Apêndice Técnico.
\end{identificacao}

\bloco{Problema e tese}
O que explica a integração financeira internacional antes de 1914 e a
convergência dos juros soberanos entre 1880 e 1913? A resposta é que a
convergência veio da queda do peso do serviço da dívida sobre a receita
pública, sustentada por crescimento econômico, e não da difusão do padrão-ouro.
A adesão ao ouro não funcionava como selo de boa conduta: era, quando muito,
seguro contra crise cambial em país endividado em moeda estrangeira, e sua
significância nas regressões anteriores vem de má especificação.

\bloco{Argumento}
O livro inteiro propõe um método indutivo para o estudo das percepções de
mercado. Em vez de projetar teorias modernas sobre dados antigos, os autores
reconstroem as ``teorias em uso'' pelos analistas de risco soberano da época,
sobretudo o Service des Études Financières do Crédit Lyonnais, e daí derivam as
variáveis que a regressão deve testar. A base de 17 países e 34 anos é montada
a partir de arquivo bancário e de fontes contemporâneas, não de retrospectivas
oficiais, porque os investidores da época corrigiam vieses que a estatística
posterior incorporou. Do exercício sai uma reinterpretação do padrão-ouro
clássico e uma tese sobre desenvolvimento: sem crescimento não há globalização
financeira.

\begin{enumerate}[leftmargin=*,nosep]
\item A integração financeira, medida por Feldstein e Horioka, oscila junto com
o prêmio de risco dos países pobres em capital (p. 17-18). Prêmio e integração
são dois lados da mesma moeda, o que justifica explicar a integração
explicando o prêmio.
\item A explicação convencional atribui a convergência ao ouro. Bordo e
Rockoff (1996) são identificados como os pioneiros da versão ``selo de
aprovação'', com redução de cerca de 50 pontos-base associada ao clube do ouro
(p. 21, nota 22, p. 60). Os autores refazem o teste da forma mais favorável
possível ao ouro, com a adesão como única variável ao lado de controles de
país, e obtêm 150 pontos-base (p. 22, nota 26, p. 60). Ainda assim a
convergência líquida da contribuição do ouro quase não difere da observada:
\chave{The gold standard explains at best only a tiny fraction of the
convergence of interest rates and thus only a fraction of financial
globalisation}{p. 22}
\item O motivo é de identificação, não de magnitude. Os anos de adesão são
definidos \emph{ex post} por estabilidade cambial, e estabilidade cambial já
significa ausência de crise, boa reputação e ambiente favorável, de modo que a
\emph{dummy} de ouro faz proxy de outra coisa (p. 24). Reformas vêm em bloco: o
Japão de 1897 adota o ouro no mesmo ano em que conclui a reforma dos direitos
de propriedade, vence a China e deposita a indenização em Londres (p. 24-25).
A seção III compara isso à estatística racialista do século XIX, que
classificava nações ``latinas'' como imprudentes e depois remanejava os casos
desviantes para salvar a classificação (p. 26).
\item A survey das fontes contemporâneas indica que o investidor monitorava o
peso da dívida, definido como serviço de juros sobre receita tributária. Nas
estimativas, um aumento de 10 pontos percentuais nesse peso eleva o juro em 70
a 80 pontos-base; renegociação de dívida acrescenta cerca de 500 pontos-base, e
a memória do calote ainda cobra 45 pontos-base dez anos depois do acordo
(p. 39). A \emph{dummy} de ouro nunca é significativa e às vezes muda de sinal
por subperíodo (p. 39).
\item O segundo modelo troca o \emph{spread} pela probabilidade de \emph{default}
e lineariza a relação. Confirma o peso da dívida, a memória do calote e a
crise política, e mantém o ouro na fronteira da não significância (p. 43-44):
\chave{Investors could see through the veil of monetary regimes}{p. 44}
\item A convergência é então explicada pela queda dos pesos da dívida
(p. 45). A decomposição mostra que o serviço bruto da dívida seguiu subindo na
maioria dos países, e que quem compensou foi o crescimento do produto, não o
equilíbrio orçamentário nem a ampliação da tributação (p. 46).
\item A conclusão reordena a hierarquia: o regime cambial importa apenas sob
``original sin'', quando o país não consegue tomar emprestado na própria moeda
e acumula exposição cambial; nesse caso a adesão do tipo \emph{currency board}
opera como seguro contra crises gêmeas, não como sinal (p. 56). A escolha do
regime cambial é consequência, e não causa, da globalização (p. 56).
\end{enumerate}

\bloco{Conceitos e definições operacionais}
Padrão-ouro internacional é arranjo informal em que cada país fixa e defende
por conta própria o valor da moeda em ouro, gerando câmbio fixo de fato entre
os aderentes (p. 11). Adesão é operacionalizada como câmbio dentro dos pontos
do ouro por pelo menos seis meses consecutivos no ano (p. 37). Peso da dívida
é serviço de juros sobre receita do governo (p. 36). O teste ``trade'' divide o
serviço da dívida pelas exportações, com zona de perigo acima de 20 por cento,
convertida do limiar de Fenn's Compendium (nota 94, p. 66); o teste ``tax''
divide pela receita pública, com limiar de 35 por cento (p. 9). Os limiares são
citação de época:
\chave{Whenever the interest service to government revenue ratio is larger
than 35 per cent, the greatest prudence is in order, although creditors can
still feel reasonably confident}{p. 41}
Reputação tem dois componentes operacionais: memória de calote, penalidade
decrescente após o acordo, e regime político, medido pela fração da população
com direito de voto (p. 36-37). Integração financeira é a correlação em corte
transversal entre poupança e investimento (p. 17). Probabilidade de
\emph{default} é derivada do próprio \emph{spread} (nota 86, p. 65).

\bloco{Evidência e método}
Base anual de 1880 a 1913 para 17 países, capital-ricos e capital-pobres,
com 480 observações nas regressões. Fontes: planilhas do Crédit Lyonnais,
Statesman's Year Book, Fenn's Compendium, The Economist e Le Rentier, mais
séries nacionais. Dois modelos com efeitos fixos de país: o modelo I explica o
diferencial de juros contra a Grã-Bretanha, o modelo II a probabilidade de
\emph{default}. O único contrafactual formal é a simulação argentina: estimado
o modelo em 1900-1913 e aplicado aos dados de 1884-1890, o teste ``trade''
reproduz o \emph{spread} observado, e o teste ``tax'' preveria até 200
pontos-base a mais em 1889 (p. 53). O deslocamento do teste ``trade'' para o
``tax'' é datado da crise Baring de 1890 (p. 52-53). No corte de 1888,
Argentina, Grécia, Itália e Portugal caem na zona de perigo dos dois testes e
Brasil e Espanha ficam na fronteira; Argentina, Brasil e Grécia pontuam pior
pelo teste ``tax'' do que pelo ``trade'', e esses três, mais Portugal, foram
os que quebraram na década de 1890. A lição que os autores atribuem ao mercado
é que um país podia ser muito aberto ao comércio, e o exemplo citado é o
Brasil, e ainda assim ter sistema fiscal profundamente deficiente (p. 52).

\bloco{O Brasil na base}
O Brasil está entre os 17 países (nota 8, p. 59). Na tabela de anos em ouro, o
país aparece com 1 apenas em 1889 e no bloco de 1906 a 1913 (Tabela DB.20,
p. 129), isto é, a \emph{dummy} brasileira de ouro coincide exatamente com a
Caixa de Conversão, que o livro nunca nomeia. O rendimento brasileiro está em
4,90 em 1905 e 4,91 em 1906, sobe a 5,31 em 1907 e volta a 4,90 em 1913,
contra 2,81, 2,86, 3,00 e 3,43 do britânico (Tabela DB.16, p. 125), ou seja, o
diferencial encolhe ao longo do período da Caixa sem queda no ano da adesão. O
Brasil
integra o grupo que entrou em crise de liquidez entre 1890 e 1898 por dívida
externa em moeda estrangeira, ao lado de Argentina, Portugal e Grécia (p. 33).
A tabela de renegociações registra problema em 1897, negociação em 1897 e
acordo em 1898 com Rothschild, sem redução de serviço (Tabela B.1, p. 99). Nos
eventos políticos constam abolição de 1888, proclamação da República em 1889,
golpe de 1891 e rebelião de 1893 (Tabela DB.22, p. 131). A franquia eleitoral
brasileira é 0 até 1893 e fica entre 1,4 e 3,3 por cento depois (Tabela DB.21,
p. 130). As fontes brasileiras são Peláez e Suzigan, Villanova Villela, IBGE,
Mitchell, o Brazilian Yearbook de 1908 e o próprio Crédit Lyonnais
(p. 99-104).

\bloco{Agentes e vertentes}
Alvos diretos: Bordo e Rockoff (1996), autores da tese do selo; Obstfeld e
Taylor, que atribuem a integração à remoção de barreiras e à estabilidade
cambial; Sussman e Yafeh no caso japonês; Clemens e Williamson; Eichengreen e
Temin para a leitura da ``mentalidade'' do ouro. Ferguson e Batley são
contestados quanto à política como fator autônomo, já que aqui a política age
sobre o crédito pela dinâmica da dívida (p. 35). Como fontes de época
aparecem Paul Leroy-Beaulieu, Robert Lucas Nash e Fenn's Compendium, o Crédit
Lyonnais, o Council of Foreign Bondholders, James de Rothschild, Laffitte,
Adolph Wagner, Baxter e Raffalovich. Filiação: Gerschenkron para o Estado
promotor do desenvolvimento, North e Weingast para instituições
representativas e crédito, Kindleberger para as manias.

\bloco{Críticas e limites}
O livro trata só do lado da oferta, isto é, da precificação pelo mercado, e
deixa fora a política de endividamento dos governos (p. 46 e p. 57). A crítica
metodológica à \emph{dummy} de ouro atinge o próprio instrumento dos autores,
que constroem uma \emph{dummy} pelo mesmo critério de estabilidade cambial e
declaram sentir por ela ``limited liability'' (p. 37). O regime político é
reduzido a uma variável, a fração de eleitores, e não há controles temporais
por falta de mudanças de regime na amostra (p. 37); a variável é insignificante
no modelo I e forte no modelo II, o que fragiliza a leitura da reputação
política. Os limiares de 20 e 35 por cento são regras de bolso de banqueiros de
Paris e Londres, e não há demonstração de que circulassem na imprensa
periférica. Para o caso brasileiro, o livro só chega até 1913, não discute a
suspensão de 1914, não nomeia agente brasileiro algum e não distingue paridade
nova de paridade antiga, distinção que em Bordo e Rockoff é a diferença entre
seguir a regra e quebrá-la.

\begin{dialogo}
\item Procurar nos editoriais e nos Documentos Parlamentares o vocabulário do
serviço da dívida contra receita e contra exportações: juros da dívida externa,
encargos, orçamento, alfândega, valor das exportações de café. Testa se a razão
pública brasileira de 1906 raciocinava em algo próximo ao teste ``tax'' que os
autores identificam como norma depois de 1890, ou se ficava no argumento
cambial do exportador.
\item Procurar se a defesa da Caixa de Conversão invocou crédito externo e
confiança de Londres, ou proteção do serviço da dívida em ouro contra a
oscilação do mil-réis. Testa a tese central do livro, de que o câmbio importava
apenas pelo canal material da dívida em moeda estrangeira e não como sinal
(p. 56).
\item Procurar menções ao funding loan de 1898 e aos Rothschild em 1906. O
livro estima que a memória do calote ainda cobrava 45 pontos-base dez anos
depois do acordo (p. 39). Testa se a imprensa tratava o crédito como estoque de
reputação a ser reconstruído, e liga a Parte II do capítulo.
\item Procurar se alguém associa regime político, representação parlamentar ou
qualidade da administração ao custo do crédito externo. Testa, num país cuja
franquia na base nunca passa de 3,3 por cento, a associação entre extensão do
voto e queda da probabilidade de \emph{default} (p. 44).
\end{dialogo}

\usonocapitulo{Parte I, como contraditório direto de Bordo e Rockoff (1996). O
livro sustenta que os fundamentos fiscais, e não a adesão ao ouro, explicavam
os \emph{spreads}, e que a \emph{dummy} de ouro só é significativa quando
monopoliza a regressão. Isso permite abrir a discussão de credibilidade sem
tratar o padrão-ouro como causa, e reposicionar a Caixa de Conversão como
instrumento de estabilização cambial de um devedor em moeda estrangeira, não
como candidatura a um clube. Serve também à Parte II, pelo registro do acordo
de 1898 com os Rothschild e pela penalidade duradoura do calote. O achado mais
útil ao capítulo é negativo: a \emph{dummy} brasileira de ouro da base é
exatamente 1906-1913, e o livro não escreve o nome da Caixa de Conversão em
nenhuma página.}
